\section{Modultests}
Die Modultests dienen der Überprüfung einzelner Hardware- und Softwarefunktionen.  
Dabei erfolgt ein Soll-Ist-Vergleich der jeweiligen Teilfunktion.  

Zunächst werden die einzelnen \gls{led}-Streifen unabhängig voneinander mit den zugehörigen Funktionen getestet.  
Hierbei wird überprüft, ob die Logik zur Ansteuerung der verschiedenen Wörter korrekt funktioniert und ob die vorgesehenen Farben sowie Helligkeiten korrekt dargestellt werden.  
Für die Hauptanzeige werden verschiedene Uhrzeiten, Helligkeiten und Farbszenarien getestet.  
Dabei wird überprüft, ob die korrekten Sätze mit den vorgesehenen Farb- und Helligkeitswerten dargestellt werden.  

Anschließend erfolgt die Prüfung der einzelnen Zusatzfunktionen.  
Hierzu werden verschiedene Werte für Wetterdaten, Temperaturen und Zeitinformationen vorgegeben.  
Dadurch kann überprüft werden, ob die korrekten Wörter und Zusatzanzeigen aktiviert werden.  
Die Funktionen zur Anzeige der Außentemperatur und der Wetterinformationen werden nacheinander geprüft.  
Zur Überprüfung der Wetterklassifikation wird kontrolliert, ob verschiedene Wettercodes und Temperaturwerte korrekt den definierten Wetter- und Temperaturgruppen zugeordnet werden.  
Nachdem dies bestätigt wurde, wird überprüft, ob die von der \gls{api} bezogenen Wetter- und Temperaturdaten die erwartete Ausgabe liefern und den realen Zuständen entsprechen.  

Nach erfolgreicher Prüfung erfolgen die Modultests der Zusatzfunktionen Innenraumtemperatur, Luftfeuchtigkeit, Tageszeit und Mondphase.  
Hierbei wird ein Fehler innerhalb der Mondphasenberechnung festgestellt.  
Die berechnete Mondphase stimmt dabei nicht mit den erwarteten Werten überein.  
Durch Anpassung des Berechnungsverfahrens kann dieser Fehler behoben werden.  

Die Modultests zeigen, dass die einzelnen Hardware- und Softwarekomponenten nach der Fehlerkorrektur zuverlässig funktionieren und die Softwaremodule unabhängig voneinander korrekt arbeiten.  